\section{Mode Memori Berpenunjuk dan Segmen Default}
\label{sec:mode-memori-segmen}

Empat mode berikutnya menyatakan alamat melalui register penunjuk,
bukan bilangan tetap. Karena isi register dapat diubah, alamatnya pun
dapat berubah saat program berjalan; kemampuan inilah yang membuat
akses berulang terhadap array dan \eng{record} menjadi mungkin
\cite{kader2014}.

\subsection{Mode Register Indirect}

Pada mode \term{register indirect}, operand ditulis di antara tanda
kurung siku berisi sebuah register. Register itu tidak menyimpan
data, melainkan \emph{offset} tempat data berada.

\begin{lstlisting}[language=MASM, numbers=none]
    MOV BX, OFFSET data1    ; BX <- offset awal array data1
    MOV AX, [BX]            ; AX <- isi memori DS:BX
    ADD BX, 2               ; maju dua byte (satu elemen word)
    MOV AX, [BX]            ; AX <- elemen kedua array
\end{lstlisting}

Hanya \reg{BX}, \reg{BP}, \reg{SI}, dan \reg{DI} yang boleh menjadi
penunjuk alamat. \reg{BX} dan \reg{BP} adalah \term{register basis}
(\eng{base register}), sedangkan \reg{SI} dan \reg{DI} adalah
\term{register indeks} (\eng{index register}).

\important{Segmen \code{[BX]} dan \code{[BP]} berbeda.} \code{[BX]}
memakai \reg{DS}, sedangkan \code{[BP]} memakai \reg{SS}; instruksi
\lstasm|MOV AX, [BP]| karena itu membaca stack, bukan segmen data.
Inilah alasan \reg{BP} dipakai pada Bab~\ref{chap:prosedur}.

\subsection{Mode Indexed dan Base-Relative}

Bentuk \code{[SI+disp]} dan \code{[DI+disp]} disebut mode
\term{indexed}; \code{[BX+disp]} dan \code{[BP+disp]} disebut mode
\term{base-relative}.

\begin{lstlisting}[language=MASM, numbers=none]
    ; --- indexed: register indeks + perpindahan ---
    MOV SI, 4               ; SI = indeks elemen ke-2 (2 x 2 byte)
    MOV AX, data1[SI]       ; AX <- memori DS:(data1 + SI)
    MOV AX, [SI+8]          ; AX <- memori DS:(SI + 8)

    ; --- base-relative: register basis + perpindahan ---
    MOV BX, OFFSET data1    ; basis = awal array
    MOV AX, [BX+4]          ; AX <- elemen ke-2
    MOV BP, SP              ; basis = puncak stack
    MOV AX, [BP+2]          ; AX <- parameter pada SS:(BP + 2)
\end{lstlisting}

Penulisan \code{data1[SI]} diterjemahkan assembler menjadi
\code{[data1 + SI]}, menyerupai penulisan array pada bahasa tingkat
tinggi. Perpindahan boleh berukuran 8 bit, 16 bit, atau nol.

\subsection{Mode Based-Indexed}

Mode \term{based-indexed} menggabungkan satu basis, satu indeks, dan
satu perpindahan sekaligus. Hanya empat kombinasi yang legal:
\code{[BX+SI+disp]}, \code{[BX+DI+disp]}, \code{[BP+SI+disp]}, dan
\code{[BP+DI+disp]}.

\begin{lstlisting}[language=MASM, numbers=none]
    MOV BX, OFFSET tabel    ; basis  : awal tabel record
    MOV SI, 6               ; indeks : record ke-3 (3 x 2 byte)
    MOV AX, [BX+SI+2]       ; AX <- medan ke-2 dari record ke-3
\end{lstlisting}

Mode ini dipakai bila alamat akhir tersusun dari dua komponen yang
bergerak sendiri-sendiri, misalnya baris dan kolom matriks, atau
nomor \eng{record} dan nomor medannya.

\subsection{Menghitung Alamat Efektif}

Prosesor mula-mula membentuk \term{alamat efektif} (\eng{effective
address}, EA), yaitu jumlah basis, indeks, dan perpindahan. Hasilnya
berupa offset 16 bit yang kembali ke nol setelah melewati batas
segmen. Alamat fisik dibentuk dengan menambahkan basis segmen
(Bab~\ref{chap:segmentasi}).

\begin{figure}[H]
    \centering
    \begin{tikzpicture}[
        komp/.style={
            draw=obeBlue, line width=0.8pt, fill=obeBlueLight,
            rounded corners=3pt, minimum width=3.9cm, minimum height=1.15cm,
            align=center, inner sep=4pt, font=\small
        },
        jumlah/.style={
            draw=obeGray, line width=0.9pt, fill=obeGrayLight,
            circle, minimum size=1.0cm, font=\large\bfseries
        },
        hasil/.style={
            draw=obeTeal, line width=1.0pt, fill=obeTealLight,
            rounded corners=3pt, minimum width=4.4cm, minimum height=1.15cm,
            align=center, inner sep=4pt, font=\small
        },
        panah/.style={-{Stealth[length=2.5mm]}, line width=0.9pt, obeBlue},
        kecil/.style={font=\scriptsize\itshape, text=obeGray}
    ]
        \node[komp] (base) at (0, 2.6)
            {\textbf{Basis} (\eng{base})\\[1pt]
             {\scriptsize \reg{BX} atau \reg{BP} atau tidak ada}};
        \node[komp] (index) at (0, 1.3)
            {\textbf{Indeks} (\eng{index})\\[1pt]
             {\scriptsize \reg{SI} atau \reg{DI} atau tidak ada}};
        \node[komp] (disp) at (0, 0)
            {\textbf{Perpindahan} (\eng{displacement})\\[1pt]
             {\scriptsize 8 bit, 16 bit, atau 0}};

        \node[jumlah] (plus) at (4.9, 1.3) {$+$};

        \node[hasil] (ea) at (9.0, 1.3)
            {\textbf{Alamat efektif (EA)}\\[1pt]
             {\scriptsize \code{EA = base + index + disp}}};

        \draw[panah] (base) -- (plus);
        \draw[panah] (index) -- (plus);
        \draw[panah] (disp) -- (plus);
        \draw[panah] (plus) -- (ea);

        \node[hasil, below=1.5cm of ea] (fisik)
            {\textbf{Alamat fisik}\\[1pt]
             {\scriptsize \code{(segmen x 16) + EA}}};
        \draw[panah] (ea) -- (fisik)
            node[midway, right=3pt, kecil] {ditambah basis segmen};
    \end{tikzpicture}
    \caption{Perhitungan alamat efektif sebagai jumlah basis, indeks,
    dan perpindahan, lalu pembentukan alamat fisiknya.}
    \label{fig:alamat-efektif}
\end{figure}

Misalkan \reg{DS} = \hex{1000}, \reg{BX} = \hex{0200}, \reg{SI} =
\hex{0006}, dan instruksi \lstasm|MOV AX, [BX+SI+2]|. Alamat
efektifnya \hex{0200} + \hex{0006} + \hex{0002} = \hex{0208}, dan
alamat fisiknya \hex{1000} $\times$ \hex{10} + \hex{0208} =
\hex{10208}. Alamat fisik inilah yang muncul pada jendela memori
\eng{debugger}.

\subsection{Kombinasi Register yang Legal dan Ilegal}

Tidak semua penulisan operand berpenunjuk dapat diterima assembler;
Tabel~\ref{tab:kombinasi-register} merangkum keduanya.

\begin{table}[H]
    \centering
    \caption{Kombinasi register penunjuk alamat pada 8086.}
    \label{tab:kombinasi-register}
    \small
    \begin{tabular}{@{}p{3.0cm}p{1.6cm}p{1.9cm}p{7.4cm}@{}}
        \toprule
        \textbf{Bentuk operand} & \textbf{Legal} &
        \textbf{Segmen} & \textbf{Alasan} \\
        \midrule
        \code{[BX]}, \code{[SI]}, \code{[DI]} & Ya & \reg{DS} &
        Tersedia \eng{encoding} untuk satu penunjuk \\
        \addlinespace
        \code{[BP]} & Ya & \reg{SS} &
        \reg{BP} adalah basis segmen stack \\
        \addlinespace
        \code{[BX+SI]}, \code{[BX+DI]} & Ya & \reg{DS} &
        Satu basis ditambah satu indeks \\
        \addlinespace
        \code{[BP+SI]}, \code{[BP+DI]} & Ya & \reg{SS} &
        Basis \reg{BP} ditambah satu indeks \\
        \addlinespace
        \code{[AX]}, \code{[CX]}, \code{[DX]} & Tidak & --- &
        Bukan basis maupun indeks; tidak ada kodenya \\
        \addlinespace
        \code{[SP]} & Tidak & --- &
        \reg{SP} bukan basis; gunakan \reg{BP} \\
        \addlinespace
        \code{[BX+BP]} & Tidak & --- &
        Dua register basis tidak dapat digabung \\
        \addlinespace
        \code{[SI+DI]} & Tidak & --- &
        Dua register indeks tidak dapat digabung \\
        \addlinespace
        \code{[BX+SI+DI]} & Tidak & --- &
        Maksimal satu basis ditambah satu indeks \\
        \bottomrule
    \end{tabular}
\end{table}

\begin{peringatanbox}
Pesan \file{Invalid use of register} pada operand seperti \code{[AX]}
bukan kesalahan penulisan semata: 8086 memang tidak memiliki
\eng{encoding}-nya, sesuai struktur \eng{ModR/M} pada
Bab~\ref{chap:format-instruksi}.
\end{peringatanbox}

\subsection{Segmen Default dan Segment Override}

Setiap operand memori dibaca melalui salah satu register segmen. Bila
tidak disebutkan, prosesor memakai segmen \eng{default} yang
bergantung pada register penunjuknya \cite{evans2020}.

\begin{figure}[H]
    \centering
    \begin{tikzpicture}[
        opd/.style={
            draw=obeBlue, line width=0.8pt, fill=obeBlueLight,
            rounded corners=3pt, minimum width=3.2cm, minimum height=0.9cm,
            align=center, font=\small\ttfamily
        },
        seg/.style={
            draw=obeTeal, line width=0.9pt, fill=obeTealLight,
            rounded corners=3pt, minimum width=1.8cm, minimum height=0.9cm,
            align=center, font=\small\bfseries
        },
        panah/.style={-{Stealth[length=2.5mm]}, line width=0.9pt, obeTeal},
        kecil/.style={font=\scriptsize\itshape, text=obeGray}
    ]
        \node[opd] (bx) at (0, 1.8) {[BX + disp]};
        \node[opd] (bp) at (0, 0.4) {[BP + disp]};
        \node[seg] (ds) at (6.2, 1.8) {DS};
        \node[seg] (ss) at (6.2, 0.4) {SS};

        \draw[panah] (bx) -- (ds) node[midway, above, kecil] {default};
        \draw[panah] (bp) -- (ss) node[midway, above, kecil] {default};

        \node[kecil, below=0.55cm of bp, align=center]
            {\eng{Segment override prefix} mengubah segmen yang
             dipakai,\\ tetapi tidak mengubah nilai offsetnya.};
    \end{tikzpicture}
    \caption{Segmen \eng{default} untuk \code{[BX+disp]} adalah
    \reg{DS}, sedangkan untuk \code{[BP+disp]} adalah \reg{SS}.}
    \label{fig:segmen-ds-ss}
\end{figure}

\begin{table}[H]
    \centering
    \caption{Segmen \eng{default} setiap bentuk operand memori 8086.}
    \label{tab:segmen-default}
    \small
    \begin{tabular}{@{}p{5.6cm}p{2.6cm}p{6.2cm}@{}}
        \toprule
        \textbf{Bentuk operand} & \textbf{Segmen} &
        \textbf{Keterangan dan override} \\
        \midrule
        \code{[disp]}, nama variabel & \reg{DS} &
        Override: \code{CS:}, \code{SS:}, atau \code{ES:} \\
        \addlinespace
        \code{[BX]}, \code{[SI]}, \code{[DI]}, \code{[BX+SI]},
        \code{[BX+disp]} & \reg{DS} &
        Penunjuk data; override \code{SS:}, \code{ES:},
        \code{CS:} \\
        \addlinespace
        \code{[BP]}, \code{[BP+disp]}, \code{[BP+SI]},
        \code{[BP+DI]} & \reg{SS} &
        Penunjuk stack; override \code{DS:}, \code{ES:},
        \code{CS:} \\
        \addlinespace
        Operand string \reg{DI} sebagai tujuan & \reg{ES} &
        Tidak dapat di-override \\
        \addlinespace
        Pengambilan instruksi (\reg{IP}) & \reg{CS} &
        Tidak dapat di-override \\
        \addlinespace
        \lstasm|PUSH| dan \lstasm|POP| & \reg{SS} &
        Selalu memakai stack; tidak dapat di-override \\
        \bottomrule
    \end{tabular}
\end{table}

Override ditulis dengan menempatkan nama register segmen di depan
tanda kurung, misalnya \lstasm|MOV AX, ES:[BX]|. Override hanya
mengganti register segmen, bukan offset, dan menambah satu byte
awalan (\eng{prefix}).
